\section{Random input domains and equivariant kernels}
\label{lifts:random-marginals}

Lemma~\ref{lifts:averaged-tuples} takes the image-law error with the
input tuple fixed and bounds its average over the tuple.
We use this information in three ways. One average selects a partition
whose complementary marginals admit a common truncation. A stronger
average controls sixteen separately truncated transition kernels. A
third argument keeps all omitted sets and uses the companion's
projection expansion for those sets.

For a probability $\mu$ on $G=\operatorname{Sym}(V)$ and a specified
set $A\subseteq V$, write $\mu_A$ for the law of the restriction
$g|_A$ under $g\sim\mu$, and write $U_A$ for the uniform law on
injections from $A$ to $V$. Choosing an ordering of $A$ identifies
these laws with the corresponding ordered image laws. A different
ordering only permutes coordinates, so it preserves their
total-variation distance. Thus the averages in
Lemma~\ref{lifts:averaged-tuples} and
Corollary~\ref{lifts:tuple-constants} also give the corresponding
averages over uniformly chosen input sets.

\subsection{Four blocks from sixteen complementary marginals}
\label{lifts:four-blocks-section}

The first application uses the full-partition case of the isotypic
estimate in Proposition~\ref{shear:prop:disjoint-transfer}, proved for
partitions in the companion \emph{From partial permutation information
to Fourier bounds} \cite[Proposition~\Lref{l:isotypic}]{partial-fourier}.
Let $V=M_1\sqcup\cdots\sqcup M_{16}$ be a partition into nonempty label
sets, and put $H_i=\operatorname{Sym}(M_i)$, fixing the complement
pointwise. If $f\ge0$ has mass $z\le1$ and satisfies
$f(gH_i)\le2/[G:H_i]$ for every $g,i$, then for every irreducible
unitary representation $\rho$ of degree $D$,
\begin{equation}
 \|\widehat f(\rho)\|_{\HS}^2
 \le32zC_2D^{-3/4}.
 \label{lifts:isotypic-input}
\end{equation}
The transform is not divided by $z$, the Hilbert--Schmidt norm uses
the ordinary trace, and the estimate includes every restriction
multiplicity. This is \eqref{shear:eq:noncovering-isotypic-general}
with $r=16$, $K=2$, and $u=2$; that equation also allows disjoint
supports that leave labels unused. For the two convolution arguments
below, the direct inverse-square
estimate \cite[Lemma~\Lref{l:degree-inverse-square}]{partial-fourier}
supplies
\begin{equation}
 C_2<\infty,\qquad
 Z_2(n)=\sum_{\substack{\rho\in\operatorname{Irr}(S_n)\\
                         D_\rho>1}}D_\rho^{-2}\longrightarrow0,
 \qquad D_\rho=\dim\rho.
 \label{lifts:inverse-square-input}
\end{equation}

\begin{proposition}\label{lifts:random-four-blocks}
As $d\to\infty$, with $n=2^d$,
\[
 \max_{g_0\in G}\|P_d^{2052d}(g_0,\cdot)-U_n\|_{\TV}
 =\|q_d^{*(2052d)}-U_n\|_{\TV}\longrightarrow0.
\]
\end{proposition}
\begin{proof}
For $d\ge4$, put $m=n/16$, $T=513d$, and $\mu=q_d^{*T}$.
Corollary~\ref{lifts:tuple-constants} gives
$\mathbb E_{|A|=n-m}\|\mu_A-U_A\|_{\TV}\le\varepsilon_n=o(1)$.
In a uniformly chosen ordered partition $(M_1,\ldots,M_{16})$
into sets of size $m$, each complement $V\setminus M_i$ is uniform.
Hence the expected sum of the sixteen complementary errors is at most
$16\varepsilon_n$, and some partition has sum at most this value.

For that partition, the restriction to $V\setminus M_i$ identifies
the left coset $gH_i$. Delete the union of all cosets whose original
$\mu$-mass exceeds twice uniform. The calculation in the proof of
Lemma~\ref{lem:trim} gives a common subprobability $f\le\mu$ of mass
$z=1-\eta$, with $\eta\le32\varepsilon_n=o(1)$, such that
$f(gH_i)\le2/[G:H_i]$ for every $g,i$.
Equation~\eqref{lifts:isotypic-input}, with $z\le1$, therefore gives
$\|\widehat f(\rho)\|_{\HS}^2\le32C_2D^{-3/4}$.
Hilbert--Schmidt submultiplicativity yields
\[
 D\|\widehat f(\rho)^4\|_{\HS}^2
 \le D\|\widehat f(\rho)\|_{\HS}^8
 \le(32C_2)^4D^{-2}.
\]
The trivial coefficient of $f^{*4}$ is $z^4$.
Lemma~\ref{lem:fair-sign} gives
$|\widehat f(\sgn)|\le\eta$, since the original block has sign mean zero.
Plancherel and \eqref{lifts:inverse-square-input} now imply
\[
 |G|\|f^{*4}-z^4U_n\|_2^2
 \le \eta^8+(32C_2)^4Z_2(n)=o(1).
\]
Cauchy--Schwarz gives $\|f^{*4}-z^4U_n\|_1=o(1)$.
Positivity gives $f^{*4}\le\mu^{*4}$, and their mass difference is
$1-z^4=o(1)$. Consequently
\[
 \|\mu^{*4}-U_n\|_{\TV}
 \le 1-z^4+\tfrac12\|f^{*4}-z^4U_n\|_1=o(1).
\]
The four independent physical blocks have total length
$4T=2052d$ and law $\mu^{*4}=q_d^{*(2052d)}$.
Equation~\eqref{eq:worst-state} gives the stated maximum over starts.
\end{proof}

\subsection{Sixteen retained kernels and the sign condition}

The second application keeps every starting deck and makes one
truncation for each omitted label set. Fix a partition
$V=M_1\sqcup\cdots\sqcup M_{16}$ with $|M_i|=m=n/16$, and put
$H_i=\operatorname{Sym}(M_i)$ and $H=\prod_iH_i$.
These groups act on decks on the right, by $g\mapsto gh$.
The quotient $X_i=G/H_i$ records the positions of all labels outside
$M_i$; it has $L=n!/m!$ states, each with $m!$ decks.

Suppose a kernel $K$ satisfies $K(gh,zh)=K(g,z)$ for $h\in H$.
Its quotient kernel is
\[
 p_i(x,y)=\sum_{z\in y}K(g,z),\qquad g\in x,\quad x,y\in X_i.
\]
Equivariance under $H_i$ makes this independent of the representative
$g$. For $\epsilon>0$, retain exactly the quotient blocks below the
specified density threshold:
\begin{equation}
 T_i(g,z)=K(g,z)\,
 \mathbf1\left\{p_i(gH_i,zH_i)\le\frac{1+\epsilon}{L}\right\}.
 \label{lifts:kernel-cutoff}
\end{equation}
Kernels act on functions by $(T_iF)(g)=\sum_zT_i(g,z)F(z)$, with
the counting inner product. In the $H_i$-isotypic decomposition,
write $T_i^{\sgn}$ for the operator on the full multiplicity space
of the sign representation, including every quotient fiber.

The sixteen-kernel theorem
\cite[Theorem~\Lref{l:equivariant-sixteen}]{partial-fourier} applies
when $K$ is doubly stochastic, $m\to\infty$, every quotient has mean
row variation at most $\delta_n\to0$, and
$\epsilon\to0$ with $\delta_n/\epsilon\to0$. Its additional input is
$\|T_i^{\sgn}\|_{\HS}^2=o(1)$ for every $i$. It then gives mean row
convergence of $K^{16}$ to uniform, accounting also for the entire
all-trivial multiplicity space. If $K$ commutes with a transitive
action on decks, all row distances are equal and the convergence is
uniform over starting decks.

\begin{theorem}\label{lifts:sixteen-theorem}
As $d\to\infty$, with $n=2^d$,
\[
 \max_{g_0\in G}\|P_d^{32784d}(g_0,\cdot)-U_n\|_{\TV}
 =\|q_d^{*(32784d)}-U_n\|_{\TV}\longrightarrow0.
\]
\end{theorem}
\begin{proof}
For sufficiently large $d$, in particular $d\ge4$, set $T=2049d$,
$\mu=q_d^{*T}$, and $K=P_d^T$.
The identity $K(g,z)=\mu(zg^{-1})$ shows that $K$ is doubly stochastic
and that it commutes with the right action of every label permutation.
A uniform source state in $X_i$ places the specified labels outside
$M_i$ at a uniform injection. Their output law is the quotient row
$p_i(x,\cdot)$, so Corollary~\ref{lifts:tuple-constants} gives
\[
 \frac1L\sum_{x\in X_i}\frac12\sum_{y\in X_i}
       |p_i(x,y)-L^{-1}|\le\delta_n,\qquad \delta_n=n^{-10}.
\]
Choose $\epsilon=n^{-1}$ in \eqref{lifts:kernel-cutoff}. A removed
quotient atom has positive excess over uniform at least
$\epsilon/(1+\epsilon)$ times its mass. Thus the mean mass removed is
\[
 \gamma_i:=\frac1{|G|}\sum_{g,z}(K(g,z)-T_i(g,z))
 \le\frac{1+\epsilon}{\epsilon}\delta_n=o(1).
\]
Each $T_i$ is row- and column-substochastic because $0\le T_i\le K$.
Moreover $H_i$ is normal in $H$, and right equivariance gives
$p_i(xh,yh)=p_i(x,y)$ for $h\in H$. The cutoff therefore preserves
the action of the whole product $H$.

We now verify the sign condition with its full multiplicity.
Fix $i$ and choose a deck $g_x$ in each fiber $x\in X_i$. For
$a,b,c\in H_i$, put $w_{xy}(c)=T_i(g_x,g_yc)$. Right equivariance gives
\[
 T_i(g_xa,g_yb)=w_{xy}(ba^{-1}),\qquad
 \sum_cw_{xy}(c)=t_i(x,y)
 :=p_i(x,y)\mathbf1\{p_i(x,y)\le(1+\epsilon)/L\}.
\]
On fiber $x$, define the normalized sign vector by
$\phi_x(g_xa)=\sgn(a)/\sqrt{m!}$, and set it to zero off that fiber.
Each fiber is a regular $H_i$-set, so these $L$ orthonormal vectors
span the full sign isotypic space of $\ell^2(G)$. In this basis the multiplicity
operator has entries
\begin{equation}
 T_i^{\sgn}(x,y)=\langle\phi_x,T_i\phi_y\rangle
 =\sum_{c\in H_i}w_{xy}(c)\sgn(c),
 \qquad |T_i^{\sgn}(x,y)|\le t_i(x,y).
 \label{lifts:sign-fiber-coefficient}
\end{equation}

Suppose the ignored positions in $x$ contain a matched pair of the
first physical shuffle, and let $s$ be its position transposition.
Then $\tau=g_x^{-1}sg_x\in H_i$ is an odd transposition and
$g_x\tau=sg_x$. In the first physical step $RS$, the switch $S$ is
uniform in $C_{e_1}$. Since $s\in C_{e_1}$, one has
$Ss\stackrel{\mathrm{law}}=S$. The first-step laws from $g_x\tau$
and $g_x$ are therefore identical, and so
$K(g_x\tau,z)=K(g_x,z)$ for every output deck $z$.
The cutoff depends only on the source and target quotient states,
so the same identity holds for $T_i$. Right equivariance now gives
\[
 w_{xy}(c)
 =T_i(g_x\tau,g_yc)
 =T_i(g_x,g_yc\tau^{-1})
 =w_{xy}(c\tau^{-1}).
\]
Since $\sgn(\tau)=-1$, the signed sum in
\eqref{lifts:sign-fiber-coefficient} is its own negative. Every
outgoing sign coefficient from this fiber is zero.

For a uniform fiber $x$, its ignored positions form a uniform
$m$-subset of $V$. The fraction of these sets containing no full
first-layer pair is
\[
 a_{n,m}=\prod_{j=0}^{m-1}\frac{n-2j}{n-j}
 \le\exp\left(-\frac{m(m-1)}{2n}\right).
\]
After $j$ positions have been sampled without a partner, exactly $j$
partners are forbidden among the $n-j$ remaining positions. The product
follows, and $\log(1-u)\le-u$ gives the exponential bound.
Only these $La_{n,m}$ fiber rows can contribute to the sign matrix.
For each such row the quotient cap gives
\[
 \sum_y|T_i^{\sgn}(x,y)|^2
 \le\sum_y t_i(x,y)^2
 \le\frac{1+\epsilon}{L}\sum_y t_i(x,y)
 \le\frac{1+\epsilon}{L}.
\]
Consequently the squared Hilbert--Schmidt norm on the entire sign
multiplicity space is
\[
 \|T_i^{\sgn}\|_{\HS}^2
 =\sum_{x,y}|T_i^{\sgn}(x,y)|^2
 \le(1+\epsilon)a_{n,m}
 \le(1+\epsilon)e^{-m(m-1)/(2n)}=o(1).
\]

Here $m=n/16\to\infty$ and $\delta_n/\epsilon=n^{-9}\to0$.
The companion theorem therefore gives mean row error $o(1)$ for
$K^{16}$. The right action of all label permutations is transitive
on decks and commutes with $K$, so every row has the same error.
Finally $K^{16}=P_d^{16T}$ and $16T=32784d$, proving the assertion.
\end{proof}

\subsection{Averaging all omitted sets}

The third application keeps the average over all omitted sets.
The omitted-set transfer
\cite[Theorem~\Lref{l:omitted-set-transfer}]{partial-fourier}
assumes, as $n\to\infty$ through multiples of $1024$, that probabilities
$\mu$ on $S_n$ satisfy
\[
 \delta_n:=\mathbb E_{|A|=n-n/1024}
       \|\mu_A-U_A\|_{\TV}\longrightarrow0.
\]
For all sufficiently large $n$, it supplies a probability $\nu$ with
$\|\nu-\mu\|_{\TV}\le16\delta_n$ and
$\|\widehat\nu(\rho)\|_{\op}\le CD^{-1/4}$ for $D=\dim\rho>1$,
where $C$ is absolute. The proof uses a finite expansion into ordered
products of projections associated with the omitted sets. The sets may overlap,
and the argument preserves the order of their projections.

\begin{proposition}\label{lifts:neumann-theorem}
As $d\to\infty$, with $n=2^d$,
\[
 \max_{g_0\in G}\|P_d^{262152d}(g_0,\cdot)-U_n\|_{\TV}
 =\|q_d^{*(262152d)}-U_n\|_{\TV}\longrightarrow0.
\]
\end{proposition}
\begin{proof}
For sufficiently large $d$, in particular $d\ge10$, set $b=1024$,
$T=(1+32b)d$, and $\mu=q_d^{*T}$.
Corollary~\ref{lifts:tuple-constants}, in the set form established at
the start of this section, gives the displayed hypothesis with
$\delta_n=o(1)$. Let $\nu$ be the resulting nearby probability.
Lemma~\ref{lem:fair-sign} gives
$|\widehat\nu(\sgn)|\le2\|\nu-\mu\|_{\TV}=o(1)$.
For eight factors, Plancherel and
$\|A\|_{\HS}\le\sqrt D\,\|A\|_{\op}$ give
\[
 4\|\nu^{*8}-U_n\|_{\TV}^2
 \le|\widehat\nu(\sgn)|^{16}
      +C^{16}Z_2(n)=o(1),
\]
using \eqref{lifts:inverse-square-input}; the nonlinear exponent is
$2-16(1/4)=-2$. Replacing the eight factors one at a time costs at
most $8\|\nu-\mu\|_{\TV}=o(1)$. Thus $\mu^{*8}$ converges to uniform.
These physical blocks have total length
$8T=8(1+32\cdot1024)d=262152d$, so
$\mu^{*8}=q_d^{*(262152d)}$. Equation~\eqref{eq:worst-state}
gives the same error from every start.
\end{proof}
